\newif\ifglobal

%%%%%%%%%%%%%%%%%%%%%%%
%%% Compile to view %%%
%%%%%%%%%%%%%%%%%%%%%%%

%%%%%%%%%%%%%%%%%%%%%%%%%%%%%%%%%%%%%%%%%%%%%%%%%%%
%%% Readme:                                     %%%
%%% https://www.overleaf.com/read/bwdjvfjksvjy/ %%%
%%%%%%%%%%%%%%%%%%%%%%%%%%%%%%%%%%%%%%%%%%%%%%%%%%%

% >>> M?: marks a module setting number or important notice

% >>> M4: Set {./relative-path} to external files for this module <<<
\def \modulefiles {./69-BUSARCH}
% To add an external file, use:
% (Replace '---' with actual filename)
% '\input{\modulefiles/tables/---.tex}' for tables
% '\includegraphics{\modulefiles/figures/---}' for figures
% '\subfile{\modulefiles/---}' for register files

% >>> M1: This module is in global project? <<<
% 'YES' - uncomment; 'NO' - comment out
\globaltrue

\ifglobal
    \documentclass[main\_\_EN.tex]{subfiles}
   
\else
    \documentclass[openany, 10pt]{book}

    % >>> M2: In ./00-shared/preamble-trm-module, also find
    %         settings for visibility of todo notes, watermark <<<
    \usepackage{preamble-trm-module}

    % >>> M3: Temporary labels in English,
    %         you can add your own labels there <<<
    \myexternaldocument{./00-shared/temp-labels-en}

\fi %\ifglobal


\setcounter{secnumdepth}{4}


% >>> M5: If this module needs any updates in
% './00-shared/preamble-trm-module.sty'
% (include additional package, etc.), contact your module PM
% or create the file './00-shared/changelog.tex' make the
% necessary updates yourself and document them in this file <<<

% Make text left-aligned and allow latex to break pages naturally, comment out for full justification
\raggedright
\raggedbottom


\begin{document}


% Sets language into which pre-defined bits of text are translated
\selectlanguage{english}
\setenglishformatting

% Do not delete this block
\ifglobal
% Add the following front matter if
% this module is outside of global project
\else
    \InputIfFileExists{readme}{}{}
    \clearpage
    \listoftodos
    \listofdonetodos
    \tableofcontents
    %\listoftables
    %\listoffigures
\fi


%%%%%%%%%%%%%%%%%%%%%%%%%
%%% TRM Module Begins %%%
%%%%%%%%%%%%%%%%%%%%%%%%%

\hypertarget{busarch}{}
\chapter{Bus Architecture}
\label{mod:busarch}

The \chipname{} SoC uses a multi-layer bus matrix as the core of its communication architecture, connecting multiple masters and slaves efficiently. This architecture supports concurrent access, effectively mitigating conflicts so that the high-performance and low-power RISC-V cores, DMA controllers, debug interfaces, and various peripherals can operate simultaneously with minimal contention. By optimizing data flow and system performance, the bus matrix ensures efficient system operation even when multiple high-speed peripherals are running at the same time.

The main system of the \chipname{} SoC consists of:

\begin{itemize}
    \item Masters:
        \begin{itemize}
            \item I-bus, D-bus and AHB bus of the high-performance (HP) 32-bit RISC-V processor 0
            \item I-bus, D-bus and AHB bus of the high-performance (HP) 32-bit RISC-V processor 1
            \item Low-power (LP) 32-bit RISC-V processor
            \item TRACE0 (Trace Interface 0)
            \item TRACE1 (Trace Interface 1)
            \item HP L2MEM MONITOR (HP L2 Memory Monitor)
            \item HP SPM MONITOR (HP Scratchpad Memory Monitor)
            \item EMAC
            \item USB OTG\_HS (USB On-The-Go High-Speed)
            \item USB OTG\_FS (USB On-The-Go Full-Speed)
            \item SD/MMC
            \item GDMA-AHB
            \item VDMA
            \item GDMA-AXI
            \item 2D-DMA
            \item H264 ENC
        \end{itemize}
    \item Slaves:
        \begin{itemize}
            \item LP ROM
            \item LP SRAM
            \item HP ROM
            \item HP SPM
            \item HP L2MEM
            \item External flash
            \item External RAM
            \item Peripherals (LP, HP and HP CPU peripherals)
            \item CSI Bridge
            \item DSI Bridge
        \end{itemize}
\end{itemize}

The bus architecture of the SoC is shown below.

\begin{figure}[H]
    \centering
    \includegraphics[width=1.0\textwidth]{\modulefiles/figures/bus_matrix.png}
    \begin{minipage}{0.9\textwidth}
    \small
    Note:\\
    1. Circles at intersections indicate that the bus (column) can access the corresponding slave device (row).\\
    2. Peripherals in the diagram include LP, HP, and HP CPU peripherals.\\
    3. CPU in the diagram represents both high-performance RISC-V processor 0 and processor 1, with each core having independent I-bus, D-bus, and AHB bus interfaces.\\
    4. The 2x MEM MONITOR in the diagram refers to HP L2MEM MONITOR and HP SPM MONITOR.\\
    \end{minipage}
    \caption{\chipname{} Bus Architecture}
    \label{fig:bus-arch}
\end{figure}

Below is the description of each bus shown in Figure \ref{fig:bus-arch} \textit{\nameref{fig:bus-arch}}:

\begin{itemize}

    \item \textbf{HP CPU I-bus}: Connects the instruction bus of the high-performance (HP) 32-bit RISC-V processor to the bus matrix for instruction fetch. It can directly access the HP SPM or indirectly access other code-containing memories through the cache, including HP ROM, HP L2MEM, as well as external flash and external RAM. In most cases, the cache is transparent to users.

    \item \textbf{HP CPU D-bus}: Connects the data bus of the high-performance (HP) 32-bit RISC-V processor to the bus matrix for literal loading and debug access. It can directly access the HP SPM or access other code- or data-containing memories (such as HP ROM, external flash, and external RAM) through the cache. It can also directly access peripheral registers, but only for HP CPU peripherals. In most cases, the cache is transparent to users.

    \item \textbf{HP CPU AHB bus}: Connects the AHB bus of the high-performance (HP) 32-bit RISC-V processor to the bus matrix for accessing data in peripherals or memory. Through the AHB interface, both LP and HP peripherals can be accessed, and direct access to HP L2MEM, HP ROM, LP SRAM, and external flash and external RAM can be configured via address mapping (refer to Chapter \ref{mod:sysmem} \textit{\nameref{mod:sysmem}} for detailed configuration of address mapping).

    \item \textbf{LP CPU bus}: Connects the low-power (LP) 32-bit RISC-V processor to the bus matrix, allowing the LP CPU to access peripherals or memory. Accessible targets include LP SRAM, LP ROM, HP L2MEM, HP ROM, external flash, external RAM, as well as LP, HP, and HP CPU peripherals.

    \item \textbf{TRACE bus}: Connects TRACE0 and TRACE1 modules to the bus matrix for debugging operations. Accessible targets include HP L2MEM and external RAM.

    \item \textbf{MEM MONITOR bus}: Connects HP L2MEM MONITOR and HP SPM MONITOR to the bus matrix for writing monitoring data from their respective monitored memory regions to memory. The two MEM MONITOR interfaces have the same access path. Accessible targets include HP L2MEM and external RAM.

    \item \textbf{EMAC bus}: Connects the Ethernet MAC controller to the bus matrix for high-speed network data transfer. Accessible targets include HP L2MEM and external RAM.

    \item \textbf{USB OTG\_HS bus}: Connects the high-speed USB OTG controller to the bus matrix for USB 2.0 high-speed data transfer. Accessible targets include HP L2MEM and external RAM.

    \item \textbf{USB OTG\_FS bus}: Connects the full-speed USB OTG controller to the bus matrix for USB 2.0 full-speed data transfer. Accessible targets include HP L2MEM and external RAM.

    \item \textbf{SD/MMC bus}: Connects the SD/MMC controller to the bus matrix for SD/MMC card data read and write operations. Accessible targets include HP L2MEM and external RAM, supporting high-speed data transfer through DMA.

    \item \textbf{GDMA bus}: The chip contains two types of GDMA controllers and one VDMA controller for high-speed data transfer.

        \begin{itemize}

            \item \textbf{GDMA-AHB}: DMA controller based on AHB bus protocol

            \item \textbf{GDMA-AXI}: DMA controller based on the AXI bus protocol

            \item \textbf{VDMA}: Video DMA controller based on the AXI bus protocol

        \end{itemize}

        Each DMA provides one connection to the bus matrix and one connection to external memory. GDMAs perform high-speed data transfers through these buses, including peripheral-to-memory, memory-to-peripheral, peripheral-to-peripheral, and memory-to-memory transfers. Accessible targets include HP L2MEM, and external flash and external RAM. For DMA accesses to external memory, it is recommended to configure longer burst lengths whenever possible to ensure optimal data throughput.

    \item \textbf{2D-DMA bus}: Connects the DMA2D graphics accelerator to the bus matrix for 2D graphics data transfer and processing. Accessible targets include HP L2MEM and external RAM.

    \item \textbf{H264 ENC bus}: Connects the H.264 video encoder to the bus matrix, containing two AXI interfaces with the same access path. Used for data access during video encoding. Accessible targets include HP L2MEM and external RAM.

\end{itemize}

\end{document}
